\section{The geometric fibre: Bianchi I from spectral data}
\label{sec:geometric_fiber}

This section attaches the emergent Bianchi-I metric $G(\mu)$ to the
spectrum of the transfer matrix $T_m$. The construction is standard
in Persistence Theory; we give only the elements needed for
Lemma~\ref{lem:G_measurable}, whose conclusion enters Step~2 of the
proof of Theorem~\ref{thm:empirical_invariance}. Readers familiar
with the holonomy chain in chapter~6 and the metric reconstruction
in chapter~12 of the monograph~\cite{PT_Monograph} may skip directly
to Section~\ref{ssec:G_measurability}.

\subsection{Holonomic angles from $T_m$}
\label{ssec:holonomic_angles}

For each prime~$p$ and each scale $\mu > 0$, the holonomy angle
$\theta_p(\mu)$ is defined by the geometric identity
\begin{equation}
  \cos \theta_p(\mu) \;:=\; 1 - \delta_p(\mu),
  \qquad
  \delta_p(\mu) \;:=\; \frac{1 - q(\mu)^p}{p},
  \label{eq:theta_p_def}
\end{equation}
where $q(\mu)$ is one of the two bifurcation values
$\qplus(\mu) = 1 - 2/\mu$ or $q_-(\mu) = e^{-1/\mu}$ singled out by
the maximum-entropy lemma L0~\cite[Ch.~3]{PT_Monograph}. As recalled
in Section~\ref{ssec:emergent_geometry}, the present paper uses the
couplage branch $q = \qplus$ throughout. By the holonomy theorem
T6~\cite[Ch.~6,~Thm.~T6]{PT_Monograph}, the squared sine of
$\theta_p$ admits the closed form
\begin{equation}
  \sinp{p}(\mu) \;=\; \delta_p(\mu)\bigl(2 - \delta_p(\mu)\bigr),
  \label{eq:sin2_T6}
\end{equation}
which is the identity used throughout the paper as a function of the
single argument $\mu$ and the choice of branch.

The interpretation in Persistence Theory is as follows. The cyclic
group $\mathbb{Z}/p\mathbb{Z}$ is the discrete analogue of a circle
of $p$ points, and $\theta_p$ is the angle accumulated by the
holonomy of the modular Fourier transform as the sieve completes one
cycle on that circle. The active primes $p \in \{3, 5, 7\}$ are the
primes whose holonomy contributes a non-trivial phase at the
canonical fixed point $\mustar = 15$; for $p \geq 11$ the
contribution is exponentially suppressed (the so-called
\emph{echo primes}, cf.~\cite[Ch.~22]{PT_Monograph}).

\subsection{Anomalous dimensions and scale factors}
\label{ssec:gamma_a}

The anomalous dimension associated with the holonomy angle
$\theta_p$ at scale~$\mu$ is the logarithmic derivative
\begin{equation}
  \gamp{p}(\mu)
  \;:=\;
  -\,\frac{\partial \ln \sinp{p}(\mu)}{\partial \ln \mu}.
  \label{eq:gamma_def}
\end{equation}
By direct differentiation of~\eqref{eq:sin2_T6} and substitution of
$q = \qplus$,
\begin{equation}
  \gamp{p}(\mu)
  \;=\;
  \frac{4\,p\,q^{p-1}(1 - \delta_p)\, /\, \mu^2}
       {(1 - q^p)(2 - \delta_p)} \cdot \mu
  \;=\;
  \frac{4\,p\,\qplus(\mu)^{p-1}\,\bigl(1 - \delta_p(\mu)\bigr)}
       {\mu\,(1 - \qplus(\mu)^p)\,(2 - \delta_p(\mu))}.
  \label{eq:gamma_explicit}
\end{equation}
The corresponding per-channel scale factor is
\begin{equation}
  a_p(\mu) \;:=\; \frac{\gamp{p}(\mu)}{\mu},
  \label{eq:a_p_def}
\end{equation}
which at the canonical fixed point $\mustar = 15$ takes the values
$a_3 \approx 0.0539$, $a_5 \approx 0.0464$,
$a_7 \approx 0.0397$ on the active primes. Both
$\gamp{p}$ and $a_p$ are real-analytic functions of $\mu$ on
$(0, \infty)$, and depend on $\mu$ through the elementary functions
$\qplus(\mu)$ and $\delta_p(\mu)$ only.

\subsection{The geometric family $\mathcal{G}$}
\label{ssec:G_family}

The Bianchi-I metric attached to the active primes
$\{3, 5, 7\}$ is the diagonal $(3+1)$-dimensional Lorentzian metric
\begin{equation}
  G(\mu) \;:=\;
  \mathrm{diag}\bigl(-1,\; a_3(\mu)^2,\; a_5(\mu)^2,\; a_7(\mu)^2\bigr),
  \qquad
  \mu \in [\mu_c, \infty),
  \label{eq:G_def_repeated}
\end{equation}
where the lower endpoint $\mu_c \approx 6.97$ is the unique solution
of $\partial^2 \ln \alpha(\mu) / \partial \mu^2 = 0$ on the
positive real axis~\cite[Ch.~12,~thm:HH]{PT_Monograph}. For
$\mu \in (0, \mu_c)$ the signature is Euclidean
($g_{00} > 0$), and for $\mu \in [\mu_c, \infty)$ it is Lorentzian
($g_{00} < 0$). Throughout this paper, we restrict to the
Lorentzian domain; the canonical fixed point of the family is the
critical point $\mustar = 15$ at which the active primes
$\{3, 5, 7\}$ are simultaneously the only primes whose anomalous
dimensions satisfy $\gamp{p}(\mu) > 1/2$, by
Theorem~T5~\cite[Ch.~8]{PT_Monograph}.

The family $\mathcal{G} := \{ G(\mu) : \mu \in [\mu_c, \infty) \}$
is the geometric object that enters the conditioning of
Theorem~\ref{thm:empirical_invariance}: the question is whether the
empirical inter-channel mutual information depends on the choice of
$G \in \mathcal{G}$.

\subsection{Measurability of $G$ with respect to the spectral
$\sigma$-algebra}
\label{ssec:G_measurability}

We now state and prove the key technical lemma. Recall that the
spectrum of $T_m$ on the active subspace
$\mathcal{S}_m = \prod_p \mathcal{S}_p$ is, by
Theorem~\ref{thm:crt_tensor} and the explicit
spectrum~\eqref{eq:Tp_spectrum} of each $T_p$, the multiplicative
orbit $\{\prod_p \lambda^{(p)}_{j_p}\}$ of the per-prime spectra
$\{1, -1/(p-2)\}$ (with the latter eigenvalue of multiplicity
$p - 2$). Denote by $\sigma(\mathrm{Spec}(T_m))$ the
$\sigma$-algebra on $\Omega$ generated by the random variables
defined as functions of this spectrum --- in particular, every
function of the holonomy angles $\theta_p$ at any scale $\mu$ is
in $\sigma(\mathrm{Spec}(T_m))$, since the angles are read off
directly from the second-largest eigenvalue of each $T_p$.

\begin{lemma}[Spectral Measurability of $G$]
\label{lem:G_measurable}
\footnote{Formally verified in Lean~4 + Mathlib as
\texttt{PT.CrtDecoupling.SpectralReduction.geometric\_spectrally\_measurable}
(finite-state form); see Section~\ref{ssec:formal_verification}.}
For every $\mu \in [\mu_c, \infty)$, the geometry $G(\mu)$ is
measurable with respect to the $\sigma$-algebra generated by the
spectrum of $T_m$:
$G(\mu) \in \sigma\bigl(\mathrm{Spec}(T_m)\bigr)$.
\end{lemma}

\begin{proof}
We exhibit each entry of $G(\mu)$ as a deterministic function of
$\mathrm{Spec}(T_m)$ and the parameter $\mu$. The (0,0) entry is the
constant $-1$, which is trivially measurable. For each active
prime~$p$ and each $\mu > 0$:

(i) The second-largest eigenvalue of $T_p$ in absolute value is
$\lambda^{(p)}_1 = -1/(p-2)$ for $p \geq 5$ (cf.\
Eq.~\eqref{eq:Tp_spectrum}) and $\lambda^{(3)}_1 = -1$. In either
case it determines $p$ uniquely (and conversely), so the map
$p \mapsto \lambda_1^{(p)}$ is a bijection between active primes
and a deterministic finite subset of $\mathrm{Spec}(T_m)$. In
particular, $p$ itself is measurable with respect to
$\sigma(\mathrm{Spec}(T_m))$.

(ii) The bifurcation value $\qplus(\mu) = 1 - 2/\mu$ is a
deterministic function of $\mu$, hence measurable with respect to
every $\sigma$-algebra. (It is a constant once $\mu$ is fixed.)

(iii) By Eq.~\eqref{eq:theta_p_def}, $\delta_p(\mu)$ is the
deterministic function $(1 - \qplus(\mu)^p)/p$ of the variables
identified in (i) and (ii), and is therefore measurable with
respect to $\sigma(\mathrm{Spec}(T_m))$. By
Eq.~\eqref{eq:sin2_T6}, $\sinp{p}(\mu)$ is the deterministic
function $\delta_p(2 - \delta_p)$ of $\delta_p$, hence also
measurable. By Eq.~\eqref{eq:gamma_explicit}, $\gamp{p}(\mu)$ is a
deterministic function of $\qplus(\mu)$ and $\delta_p(\mu)$,
hence measurable.

(iv) By Eq.~\eqref{eq:a_p_def}, $a_p(\mu) = \gamp{p}(\mu)/\mu$ is
a deterministic function of $\gamp{p}(\mu)$ and $\mu$, hence
measurable; its square $a_p(\mu)^2$ likewise.

Combining (i)--(iv) for each active prime $p \in \{3, 5, 7\}$, the
diagonal entries of $G(\mu)$ are all measurable with respect to
$\sigma(\mathrm{Spec}(T_m))$. The diagonal matrix $G(\mu)$
itself, being a deterministic function of these entries, is
therefore measurable with respect to the same $\sigma$-algebra.
\end{proof}

\begin{remark}[On the meaning of measurability here]
\label{rmk:measurable_meaning}
Lemma~\ref{lem:G_measurable} should be read as the statement that
$G(\mu)$ contains no information beyond the spectral data of
$T_m$ and the (deterministic) choice of $\mu$. There is no
``trajectory dependence'' --- no functional of the actual prime
gap sequence $\{g_n\}$ enters the construction of $G(\mu)$ once
the spectrum of $T_m$ is fixed. This is the property used in
Step~2 of the proof of Theorem~\ref{thm:empirical_invariance}, to
embed $G(\mu)$ into the shift-invariant $\sigma$-algebra
$\mathcal{I}$ on the path space~$\Omega$.
\end{remark}
